% ===================================================================
%  CADRO N.º 8 — A colleita. — A caza, a vendima, a pesca.
%  Traduction 2026 d'apres texto/io, controlee sur texto/fr, texto/es
%  et texto/pt. Consigne : voir texto/gl/10-cadro-01.tex.
%
%  L'ECART DECLARE DU BLOC c2-03-1, ET C'EST LE DEUXIEME DES TROIS.
%  L'ido y oublie le renvoi « (13) » du gobelet, que le francais
%  donne ; la colonne le rend, avec les onze qui la precedent, et
%  renvoji.py le passe sans rien dire.
%
%  LA VENDANGE EST LE SEUL LIEU DU VOLUME OU LE GALICIEN EST CHEZ LUI
%  PLUS QUE SES DEUX VOISINES, et cela se voit au nombre de mots
%  qu'il a pour la meme chose. « Bagoada » le jus qui coule de la
%  presse, « bica » la rigole, « lagar » le pressoir et le local qui
%  le contient, « pipa », « cuba », « bocoi » et « barril » pour
%  quatre contenants que le francais nomme tous « tonneau ». Une
%  langue a beaucoup de mots la ou elle a beaucoup a distinguer.
%
%  ET CE N'EST PAS UN ORNEMENT : c'est la meme regle vue une
%  troisieme fois. Le lieu de l'apprentissage decide du mot (tableau
%  10 irlandais), sa date decide de sa forme (tableau 14), son pays
%  decide de son origine (tableau 15) — et l'IMPORTANCE de la chose
%  decide de leur NOMBRE. La ou l'on vendange depuis toujours, on
%  distingue ; la ou l'on ne vendange pas, un mot suffit.
% ===================================================================

%%K t08-noto-1 noto
\VUnotes{143.2mm}{%
\textit{(*) Porque esa palabra non é precisa: trátase da colleita do liño, da
vendima, da colleita das mazás? Se cadra sería bo empregar «} \VUgras{mesono}
\textit{», proposta en Mondo XIV, p. 242. Pero ata agora esa raíz nova non foi
adoptada pola Academia e agarda polas discusións segundo o regulamento idista de
costume.}%
}

%%K t08-apar-1 apar
\VUsaut{5.0mm}
\VUtitre{15.0pt}{60}{CADRO N.º 8}
\VUsaut{3.0mm}
\VUcentre{13.2pt}{90}{\textit{Primeira escena.}}
\VUsaut{3.0mm}
\VUpk{10.2pt}{40}{A colleita (*).}
\VUsaut{4.0mm}
\VUpk{10.2pt}{40}{Os aspectos do campo.}
\VUsaut{3.0mm}

%%K t08-c1-01-1 p
1. --- Co \VUgras{verán} mudou unha vez máis o aspecto do campo. A semente confiada ao
rego xermolou; unha plantiña verde saíu da terra e botou espiga. O gran foise
formando, despois madureceu, e desde agora só cómpre recollelo.

%%K t08-c1-02-1 p
2. --- As \VUgras{flores do campo} están abertas. \VUgras{Acianos}
\textsuperscript{(1)}, \VUgras{papoulas} \textsuperscript{(2)} e \VUgras{dedaleiras}
\textsuperscript{(3)} mesturan co matiz uniforme dos trigais o contraste ledo das súas
\VUgras{corolas} frescas.

%%K t08-c1-03-1 p
3. --- As árbores froiteiras cubríronse de follas; a froita madureceu. A pequena
\VUgras{cerdeira} \textsuperscript{(4)} está chea de belas \VUgras{cereixas}
\textsuperscript{(5)} que os nenos colleitan e comen con moito pracer.

%%K t08-c1-04-1 p
4. --- Desde agora o gando pode pasar todo o día ao aire libre.

%%K t08-c1-04-2 p
Os \VUgras{años} \textsuperscript{(6)} medraron e fixéronse belas \VUgras{ovellas}
\textsuperscript{(7)} e belos \VUgras{carneiros} \textsuperscript{(8)} de mesta la.
Pacen tranquilos a \VUgras{herba} do \VUgras{pasteiro} \textsuperscript{(9)}
salferido de \VUgras{margaridas.}

%%K t08-c1-04-3 p
Axiña se comezará a tosquiar \VUgras{esas} bestas para ter a súa \VUgras{la,} coa que
se fabrica o pano da nosa roupa.

%%K t08-tit-2 sub
\VUsaut{5.0mm}
\VUpk{10.2pt}{40}{O pastor.}
\VUsaut{3.4mm}

%%K t08-c1-05-1 p
5. --- O \VUgras{pastor} \textsuperscript{(10)} semella non lles prestar moita
atención. Só se preocupa da tormenta que pasa no horizonte, e o \VUgras{gardián das
ovellas} \textsuperscript{(13)}, que achega unha \VUgras{cantimplora}
\textsuperscript{(14)} aos beizos, semella aínda máis despreocupado.

%%K t08-c1-06-1 p
6. --- A calor é opresiva, pois o \VUgras{can} \textsuperscript{(15)} vén tamén el
matar a sede; lambe a auga fresca do \VUgras{regato} \textsuperscript{(16)} e asusta
unha pobre \VUgras{ra} \textsuperscript{(17)} que se acochara nas herbas da ribeira.
Esta salta á auga clara onde florecen os \VUgras{nenúfares} \textsuperscript{(18)}.

%%K t08-c1-07-1 p
7. --- Unha \VUgras{lavandeira} \textsuperscript{(19)} báñase preto da
\VUgras{fontela} \textsuperscript{(20)} que abrolla entre dous penedos e agóchase
baixo \VUgras{os arbustos espiñentos} \textsuperscript{(21)} e as \VUgras{silveiras
de estripo} \textsuperscript{(22)}.

%%K t08-tit-3 sub
\VUsaut{5.0mm}
\VUpk{10.2pt}{40}{A colleita.}
\VUsaut{3.4mm}

%%K t08-c1-08-1 p
8. --- Para os nenos do campo, a colleita do trigo é unha época moi divertida.
Ledos, \VUgras{dan reviravoltas} \textsuperscript{(24)} sobre os \VUgras{montóns
de palla} \textsuperscript{(23)}, axudan os \VUgras{segadores} \textsuperscript{(26)} e, mentres
estes segan o \VUgras{trigal} \textsuperscript{(27)} coas súas \VUgras{fouces}
\textsuperscript{(25)} ben afiadas coa \VUgras{pedra de afiar} \textsuperscript{(30)},
eles recollen o trigo e átano en \VUgras{gavelas} \textsuperscript{(31)} ou en
\VUgras{feixes} \textsuperscript{(32)}.

%%K t08-c1-09-1 p
9. --- Ás veces permítenlle aos maiores de entre eles cortaren cunha \VUgras{fouciña}
\textsuperscript{(12)} as \VUgras{canas douradas} \textsuperscript{(28)} que levan as
\VUgras{espigas} \textsuperscript{(29)} cheas de gran; e os pequenos, vixiando o
\VUgras{gando} \textsuperscript{(33)} que pace no \VUgras{pasteiro}
\textsuperscript{(34)} baixo a sombra da grande \VUgras{tileira}
\textsuperscript{(35)}, apañan coa súa \VUgras{rede} \textsuperscript{(36)} as belas
\VUgras{bolboretas.}

%%K t08-c1-09-2 p
Algúns mesmo soben ao \VUgras{carriño} \textsuperscript{(37)} para arranxaren as
gavelas que lles pasan cunha \VUgras{forcada} \textsuperscript{(38)}.

%%K t08-c1-10-1 p
10. --- En toda a \VUgras{chaira} \textsuperscript{(39)}, onde botan a voar as
\VUgras{lavercas} \textsuperscript{(41)}, reina a animación. Todos están ocupados coa
colleita. Pero, mentres os pobres seguen a usar os vellos instrumentos,
\VUgras{fouces, fouciñas} e \VUgras{mallos} \textsuperscript{(40)}, os propietarios
ricos mandan segar o seu trigo ou o seu \VUgras{centeo} \textsuperscript{(43)} cunha
\VUgras{segadora} \textsuperscript{(42)}. Despois, sen precisaren levar as gavelas á
palleira, fanse grandes \VUgras{medas} \textsuperscript{(44)}.

%%K t08-c1-11-1 p
11. --- Entón tráese unha \VUgras{malladora} \textsuperscript{(47)} que unha
\VUgras{locomóbil} \textsuperscript{(45)} move cunha \VUgras{correa de transmisión}
\textsuperscript{(46)}, e así o gran sepárase da \VUgras{palla} sen deixar a leira
onde foi sementado.

%%K t08-tit-4 sub
\VUsaut{5.0mm}
\VUpk{10.2pt}{40}{A tormenta.}
\VUsaut{3.4mm}

%%K t08-c1-12-1 p
12. --- Moitas veces hai \VUgras{treboadas} \textsuperscript{(53)} no tempo da
colleita. Primeiro óense de lonxe os tronidos do \VUgras{trono}; un \VUgras{vendaval
do oeste} \textsuperscript{(54)} comeza a soprar e, arquexando, dobra as árbores máis
grosas do \VUgras{bosque} \textsuperscript{(49)}. \VUgras{Nubarróns}
\textsuperscript{(56)} negros asaltan o ceo; atravésanos o zigzag deslumeante do
\VUgras{lóstrego} \textsuperscript{(55)}. A \VUgras{chuvia} e ás veces a
\VUgras{sarabia} comezan a caer, esmagando as colleitas listas para seren segadas.

%%K t08-c1-13-1 p
13. --- Moitas veces o lóstrego prende lume nun edificio. Así o ano pasado o tellado
do \VUgras{muíño de vento} \textsuperscript{(50)} foi reducido a cinzas e dúas das
súas \VUgras{aspas} \textsuperscript{(51)} quedaron esnaquizadas.

%%K t08-c1-14-1 p
14. --- Despois dalgunhas horas, a calma volve nacer. Un \VUgras{arco da vella}
\textsuperscript{(52)} brillante aparece no horizonte, e a natureza retoma a súa vida
interrompida por uns intres. Os campesiños, que se refuxiaran nas \VUgras{cabanas}
\textsuperscript{(48)}, volven ao traballo e esfórzanse con coraxe por amañaren os
estragos que acaban de sufrir.

%%K t08-tit-5 sub
\VUsaut{6.0mm}
\VUcentre{13.2pt}{90}{\textit{Segunda escena.}}
\VUsaut{3.6mm}

%%K t08-tit-6 sub
\VUpk{10.2pt}{40}{A caza. --- A vendima.}
\VUsaut{1.4mm}
\VUpk{10.2pt}{40}{A pesca.}
\VUsaut{4.0mm}

%%K t08-c2-01-1 p
1. --- Arredor da fin de \VUgras{agosto} ou de mediados de \VUgras{setembro}, segundo
as rexións, ábrese a caza. Apenas os primeiros raios do sol fixeron saír os
\VUgras{lagartos} \textsuperscript{(2)} do seu burato, xa o \VUgras{cazador}
\textsuperscript{(5)} bota a camiñar cos seus \VUgras{cans} \textsuperscript{(4)}.

%%K t08-c2-02-1 p
2. --- Puxo o seu sólido \VUgras{traxe de caza} \textsuperscript{(9)} de pano groso,
calzou unhas \VUgras{polainas} de coiro, coas que conseguirá camiñar na herba húmida
onde se agochan os \VUgras{sapos} \textsuperscript{(1)}; colleu a súa
\VUgras{escopeta} \textsuperscript{(6)} e o seu \VUgras{zurrón}
\textsuperscript{(10)}, e velaquí que marcha.

%%K t08-c2-03-1 p
3. --- O fervor da persecución lévao ao medio dun viñedo onde a vendima está moi
animada. Os fillos do viticultor, que de ordinario vixían as cabras no camiño, hoxe
déixanas pacer á toa e, despois de ataren a un poste un vello \VUgras{castrón}
\textsuperscript{(3)} que non deixaría de aproveitar a súa liberdade para fuxir, tamén
eles se dan a súa parte no pracer. Un colle un acio de uvas dun \VUgras{cesto da
vendima} \textsuperscript{(12)} e diviértese facendo escorrer o \VUgras{zume}
\textsuperscript{(14)} nun \VUgras{copo} \textsuperscript{(13)} para fabricar mosto
(viño sen fermentar). Outro peitéase cunha \VUgras{coroa de follas}
\textsuperscript{(15)} que acaba de tecer. Preto deles, uns \VUgras{pardais}
\textsuperscript{(16)} sen vergoña veñen mesmo diante do lagar rapinar as uvas.

%%K t08-c2-04-1 p
4. --- Mentres os nenos se divirten, os homes traballan. Os \VUgras{vendimadores}
\textsuperscript{(18)} levan as uvas á adega en \VUgras{cestos de costas}
\textsuperscript{(17)}; un \VUgras{tanoeiro} \textsuperscript{(19)} amaña as
\VUgras{pipas} \textsuperscript{(20)}, comproba se as \VUgras{aduelas}
\textsuperscript{(22)} e os \VUgras{fondos} \textsuperscript{(21)} están ben, e manda
substituír os \VUgras{arcos} \textsuperscript{(23)} rotos. Traballou tamén as
\VUgras{cubas} \textsuperscript{(25)} nas que se botan as \VUgras{uvas}
\textsuperscript{(26)} para as facer fermentar.

%%K t08-c2-05-1 p
5. --- Cando a fermentación remata, tírase o viño e ponse o bagazo no \VUgras{lagar}
\textsuperscript{(28)} e, apertando pouco a pouco co \VUgras{fuso}
\textsuperscript{(29)}, escóase o zume, que corre nunha \VUgras{cubeta}
\textsuperscript{(32)}, e obtense aínda \VUgras{viño} \textsuperscript{(31)}.

%%K t08-tit-8 sub
\VUsaut{6.0mm}
\VUpk{10.2pt}{40}{Cervexa e sidra.}
\VUsaut{3.4mm}

%%K t08-c2-06-1 p
6. --- Na parede da \VUgras{adega} \textsuperscript{(27)} trepa unha póla de
\VUgras{lúpulo} \textsuperscript{(30)}. Alí non é máis ca unha planta de adorno, pero
nalgúns países, onde a vide é moi rara, o lúpulo cultívase para fabricar a
\VUgras{cervexa.}

%%K t08-c2-07-1 p
7. --- A pequena \VUgras{maceira} \textsuperscript{(33)} está cargada de mazás que
darán, cando estean maduras, unha excelente \VUgras{sidra.} Na súa \VUgras{gaiola}
\textsuperscript{(34)}, o \VUgras{canario} \textsuperscript{(35)} e o
\VUgras{xílgaro} \textsuperscript{(36)} miran con ollos de envexa os pardais que
brincan en liberdade.

%%K t08-c2-08-1 p
8. --- Ata onde chega a vista, na costa dos outeiros, están aliñadas as
\VUgras{estacas da vide} \textsuperscript{(40)} que sosteñen as fermosísimas
\VUgras{cepas} \textsuperscript{(39)}, cargadas de \VUgras{acios} pesados.

%%K t08-c2-08-2 p
Durante todo o ano o \VUgras{viticultor} \textsuperscript{(42)} traballou para coidar
o seu \VUgras{viñedo} \textsuperscript{(38)}, e agora é recompensado pola súa
fatiga cunha bela colleita. As \VUgras{vendimadoras} \textsuperscript{(41)} enchen as
cubas postas no carriño tirado por \VUgras{mulas} \textsuperscript{(43)}, e todos
están ledos.

%%K t08-tit-9 sub
\VUsaut{5.0mm}
\VUpk{10.2pt}{40}{A pesca.}
\VUsaut{3.4mm}

%%K t08-c2-09-1 p
9. --- O mesmo acontece cos \VUgras{pescadores de cana} \textsuperscript{(44)} que,
desde a mañá, viñeron instalarse á beira da auga coas súas \VUgras{canas de pescar}
\textsuperscript{(45)}.

%%K t08-c2-09-2 p
O \VUgras{peixe} \textsuperscript{(47)} morde o \VUgras{anzol} axiña que eles
mergullan o seu \VUgras{fío} \textsuperscript{(46)} na auga, e apenas tiveron tempo de
renovar a \VUgras{engada} cando unha nova vítima se debate na punta do fío.

%%K t08-c2-09-3 p
Aínda así, o \VUgras{pescador de rede} \textsuperscript{(48)} que, desde a súa
\VUgras{barca} \textsuperscript{(49)}, bota a \VUgras{rede} no medio do \VUgras{río}
\textsuperscript{(50)} colle máis peixes.

%%K t08-c2-10-1 p
10. --- O outono é a estación dos bos paseos: a pé, \VUgras{a cabalo}
\textsuperscript{(52)}, \VUgras{en bicicleta} \textsuperscript{(53)}, \VUgras{en
automóbil} \textsuperscript{(54)}. Os \VUgras{camiños} \textsuperscript{(51)} están
limpos e aprovéitanse os últimos días bos, pois velaquí que xa as \VUgras{andoriñas}
\textsuperscript{(56)} se xuntan nos tellados para botaren a voar a todo bater cara a
países máis quentes. A follaxe da vella \VUgras{nogueira} \textsuperscript{(57)}
amarelece, as \VUgras{noces} caen, e as \VUgras{árbores} altas dispostas en grupo coma
un \VUgras{ramallo} \textsuperscript{(58)} sobre o \VUgras{alto}
\textsuperscript{(59)} axiña quedarán sen follas.

%%K t08-c2-11-1 p
11. --- O \VUgras{sol} \textsuperscript{(60)} érguese máis tarde, os días minguan,
unha frescura ben picante comeza a sentirse á \VUgras{mañá,} e velaquí que xa bandos de
\VUgras{arceas} \textsuperscript{(61)} deixan os nosos climas pouco hospitalarios.
